How\'s the probability of default defined? 
How can the value of equity be found?

The probability of default (PD) is defined as
\( \text{PD}^P_t = \mathbb{P}[A_T < K|\mathcal{F}_t] = \Phi(-d_2) \)
and \( d_2 = \frac{\log{\frac{A_t}{K}} + (\mu - \frac{\sigma^2}{2}(T - t)}{\sigma\sqrt{T - t}} \)
(which is modified to \( \mu = r \) under the risk neutral measure). 

The value of equity is 
\( E_t = A_t \Phi(d_1) - e^{-r(T-t)}K\Phi(d_2) \)
where \( d_1 = d_2 + \sigma\sqrt{T - t} \) and we use the risk neutral measure.
To calibrate the model the unknown Value of assets and volatility is required.
With the help of the volatility of the equity (which is publicly traded),
one can simultaneously solve the defining equation for \( E_t \) and
\( \sigma_E E_t = \frac{\partial E_t}{\partial A_t}\sigma A_t = \Phi(d_1)\sigma A_t \).

The difference between this discount rate and the risk-free rate
is the Yield Spread

\( D_t = e^{-(r + YS_t)(T - t)}K \)
\( YS_t = \frac{1}{T - t}\log{\frac{K e^{-r(T-t)}}{D_t}} \)


Solve a discrete example for finding the probability of default 
and the credit/yield spread


INSERT IMAGE 13, 05b

Bonds A and B have face value of \( 1 \), but can default. What is the PD
for Bond A? What is the PD for Bond B?

Risk free rate is \(\log \left( \frac{1}{0.4 + 0.3 + 0.2} \right) = 10.536\%\)
Bond A and B both have a PD of \(1/3\)
Price of Bond A: \(0.3 + 0.2 = 0.5\)
\[\Rightarrow \$0.5 = e^{-(10.536\% + Y_SA) \cdot 1 \, \text{year}} \cdot \$1\]
\[\Rightarrow Y_SA = -( \log(0.5) + 0.10536) = 58.779\%\]
Price of Bond B: \(0.4 + 0.3 = 0.7\)
\[\Rightarrow \$0.7 = e^{-(10.536\% + Y_SB) \cdot 1 \, \text{year}} \cdot \$1\]
\[\Rightarrow Y_SB = -( \log(0.7) + 0.10536) = 25.131\%\]
Bond B defaults in a good state with low marginal utility. 
Bond A defaults in a state where marginal utility is high, i.e. where the default hurts most. 
This leads to differential pricing and thus yield spreads.
In particular: pricing solely based on credit ratings is senseless!



Model incentive compatible contracts to reduce the issue
of asymmetric information in financing

Use non-monetary/non-pecuniary penalty
e.g. 
- Borrower‘s time spent on bankruptcy proceedings
Why non-monetary?

▪ Because it does not increase financier‘s utility ex post
▪ Avoid negative consumption for borrower

Important: penalty must not be too high (valuable project won‘t
be executed then) and not too low

Let the optimal penalty function be \( \phi^\ast \), the repayment \( z(y) \),
\( R \) the accumulated (not discounted) payback
then
\( \phi^\ast(z(y)) = \max\{R - z(y), 0\} \)
Financiers want to have repayment of \( I \), on average. That is
choose \( R \) such that
\( \mathbb{E}[z(y)] = I \iff \mathbb{P}[y < R]\mathbb{E}[y|y < R] + \mathbb{P}[y \geq R]R = I \)


Show how monitoring implies the existence of financial intermediators

Monitoring is again a mechanism to reduce assymmetric information

Let \( m \) be the number monitors, each with cost \( c \)
\( n \) entrepreneurs/capital seekers with i.i.d. cash flows (the independent part is somewhat problematic)

Then, financial intermediation (=delegated monitoring)
dominates direct lending if number of projects, \( n \), is large
enough.

Proof.

The total cost of direct lending:
\( n m c  \)
Total cost delegated 
\( nc + \mathbb{E}[\phi^\ast(z(y_1, \dots, y_n)] = nc \mathbb{E}[\max{n\cdot R^n_{\text{Deposit}} - \sum_i^n y_i - c, 0}] \)
Dividing by \( n \) and \( n \to \infty \) yields
\( mc > c + \mathbb{E}[\max{I + c - \frac{1}{n} \sum_i^n y_i, 0}] \to c \)


There is a tendency for banks to grow big (i.e. lend to many
entrepreneurs)
▪ Banks (ideally) turn risky loans into much less risky deposits.
▪ If perfect diversification is not possible (small 𝑛), reduction of
risk in deposits is possible:
▪ Bank uses equity as buffer identical cash flow distributions
Using a central entity for monitoring introduces the issue of conspiracy
with hazarduous borrowers.


Prove the following statement (which is a step in determining the fair price for american puts): 
Set $L^\ast = \frac{\lambda_{-}}{\lambda_{-} -1}K$ and $\frac{1}{\sigma^2}(\lambda_{-} = -(r - \delta - \frac{\sigma^2}{2}) - \sqrt{(r - \delta - \frac{\sigma^2}{2})^2 + 4r\frac{\sigma^2}{2}})$ and $ f_L(x) = \mathbb{E}[e^{-r\tau_L}(K - S_{\tau_L})_{+}|S_0 = x]$ where $\tau_L = \inf \{u | S_u \leq L\}$. As usual $K$ is the strike, $S_t$ the price of an asset (which pays dividends $\delta$) and $r$ the risk free interest rate etc. Show with the help of Ito's Lemma that $t \mapsto e^{-rt}f_{L^\ast}(S_t)$ is a supermartingale.
